\documentclass[UTF8,a4paper,12pt]{ctexart}

\usepackage[margin=2.5cm]{geometry}
\usepackage{array}
\usepackage{booktabs}
\usepackage{caption}
\usepackage{graphicx}
\usepackage{hyperref}
\usepackage{longtable}
\usepackage{setspace}
\usepackage{tabularx}
\usepackage{enumitem}

\hypersetup{
  colorlinks=true,
  linkcolor=black,
  citecolor=black,
  urlcolor=blue
}

\setlength{\parindent}{2em}
\setlength{\parskip}{0.4em}
\setstretch{1.25}
\renewcommand{\arraystretch}{1.15}
\captionsetup{labelfont=normalfont,textfont=normalfont}
\ctexset{
  section = {format = \Large\normalfont\raggedright},
  subsection = {format = \large\normalfont\raggedright}
}
\renewcommand{\textbf}[1]{#1}

\newcolumntype{L}[1]{>{\raggedright\arraybackslash}p{#1}}
\newcolumntype{Y}{>{\raggedright\arraybackslash}X}

\begin{document}

\begin{center}
{\Large 医疗人工智能的发展图景：\\从范式演进到大模型临床部署\par}
\vspace{0.8em}
匿名作者\par
\vspace{0.4em}
2026年5月
\end{center}

\vspace{1em}
\begin{center}
摘要
\end{center}
\noindent
过去四年，医疗人工智能经历了三个层面的发展：训练范式上从任务专用监督学习走向通用基础模型，输出形式上从判别式判断走向生成式交互，模态范围上从单模态走向多模态融合。以 GPT-4、PaLM 2、Gemini 与 LLaMA 系列为代表的大语言与多模态模型，将医学问答、影像报告生成、临床决策支持与患者沟通纳入统一的生成式框架之下，并持续刷新 USMLE、MultiMedQA、MultiMedBench、MedHELM 等基准表现。本文以医疗人工智能的发展图景为主线，沿范式演进、医疗大语言模型、多模态医疗大模型与专科基础模型四条线索梳理当下医疗 AI 的能力基础（第2节至第5节），并在第6节讨论临床证据、评估方法与治理框架三方面目前仍未得到解决的若干具体问题。医疗 AI 的能力扩张由数据规模、训练目标与结构先验共同驱动，走出通用化与专科化彼此互补的技术路径。多种方法的层化叠加，构成了医疗人工智能当下能力的来源，也为后续讨论可信、可部署、可监管的实现路径奠定了基础。


\noindent 关键词：医疗人工智能；医疗大语言模型；多模态医疗大模型；专科基础模型；临床部署；可信治理

\section{引言}

过去四年，医疗人工智能经历了爆发式的快速演进。自 ChatGPT 公开发布以来，在不到三十六个月的时间里，面向医学问答的 Med-PaLM 将美国执业医师考试题集 MedQA 的准确率由 67.6\% 提升至 86.5\%[8,13]，Med-Gemini 进一步提升至 91.1\%，在修正评测题歧义后甚至达到 91.8--92.9\%[7]。同时，多模态医疗大模型也从初步接入图像，发展为能够同时处理视觉、语言、表格与基因数据的模型体系；AMIE 在 OSCE 风格诊断对话中于 30/32 个专家维度上显著优于全科医生[17]，PreA 则在两家三甲医院 2,069 名患者参与的随机对照试验中将问诊时长降低近三成，使跨科室协调度提升超过一倍[20]。

本文围绕医疗人工智能的发展图景展开，共分七节。第2节梳理 2015--2026 年医疗 AI 在训练目标、模型规模与模态广度三条主轴上的演进；第3节对比英文与中文医疗大语言模型的能力及其评估体系的变化；第4节梳理多模态医疗大模型的架构谱系；第5节论证专科基础模型在结构化电子健康记录（electronic health record, EHR）时序、影像、组学等任务上仍具有不可替代的位置；第6节讨论临床证据、评估方法与治理框架三方面尚未解决的若干问题；第7节总结全文观点。

在术语口径上，本文将医疗大语言模型（medical LLM）界定为以自然语言为主输入输出、参数规模在十亿量级以上、并在医学语料上做过领域适配（继续预训练、指令微调或强化学习）的语言模型；将多模态医疗大模型（medical LMM / VLM）界定为除文本外还能接受影像、生物信号、结构化 EHR、组学等至少一种模态作为输入的大模型；将专科基础模型界定为在单一医学子领域（如眼底、病理、CT、EHR 时序）上经自监督预训练，并能在该子领域多类下游任务上展现迁移能力的模型。医疗人工智能作为更广义的术语，涵盖上述三类以及传统任务专用监督模型。基准成绩在本文中专指静态问答或分类基准上的指标（如 MedQA 准确率），临床效用则指在真实临床工作流中以临床终点（如再入院率、误诊率、问诊时长）度量的指标。

\section{范式演进回顾：从任务专用监督到生成式智能体}

医疗人工智能在 2015--2026 年的演进可划分为三个相互衔接阶段。训练目标经历了由标签驱动进入语言驱动，再由语言驱动进入工具与反馈驱动的两次跃迁；每次跃迁都为模型拓展了新的能力边界。

\begin{figure}[htbp]
  \centering
  \fbox{\parbox{0.86\textwidth}{
    \centering
    图 2.1 占位图：三阶段训练目标演化\\[0.5em]
    阶段一为任务专用监督学习，阶段二为基础模型与多模态融合，阶段三为生成式 AI 与智能体。
  }}
  \caption{医疗 AI 三阶段训练目标演化示意。}
  \label{fig:paradigm-evolution}
\end{figure}

\noindent\textit{注：} 对比学习成为阶段二多模态对齐组件，BERT 风格双向编码器在阶段三中作为检索召回端或专科编码器继续发挥作用。

\subsection{阶段一：任务专用监督学习（约 2015--2020）}

这一阶段主要依托高质量标注数据，使用 CNN 与 RNN 进行端到端监督训练。Gulshan 等以 12.8 万张视网膜图像训练 InceptionV3，在糖尿病视网膜病变筛查上达到与眼科医师相当的灵敏度与特异性；Esteva 等以 12.9 万张皮肤镜图像训练 GoogLeNet 完成 2,032 种皮肤病变的细粒度分类；McKinney 等以英国与美国两国筛查数据训练乳腺癌检测模型，在前瞻验证中减少假阳性 5.7\%、假阴性 9.4\%[3]。截至目前，美国 FDA 已批准超过 500 项 AI/ML 医疗器械软件，其中绝大多数都属于这一阶段的产物。

尽管后续范式迅速迭代，这一阶段的许多方法仍沿用至今。\textbf{对比学习}自 CLIP、ConVIRT 起逐渐取代纯监督，把影像--报告对当作弱监督信号使用，为阶段二的多模态对齐组件奠定了基础；BioBERT、ClinicalBERT 则将 \textbf{BERT 风格的双向编码器} 适配到医学文本，成为 GatorTron、Med-BERT 等后续模型的技术基线[9]；以 \textbf{固定测试集上报告 AUC} 为核心的回顾性验证范式也在这一时期被确立，并在后续多年中保持为医疗 AI 论文默认的评测协议。

\subsection{阶段二：基础模型与多模态融合（2020--2023）}

这一时期的研究重心由单一任务的端到端训练，转向先通过自监督预训练构建通用底座，再迁移至具体任务。Transformer、AlphaFold、CLIP 与 GPT-3 共同构成了这一阶段的基础设施：Transformer 把序列建模统一成可扩展至跨模态的通用架构；AlphaFold 证明了大规模自监督在生物分子结构上的强迁移能力；CLIP 把图像与文本的对比学习推升为可扩展的视觉表示学习目标；GPT-3 则首次展示出不更新权重也能通过提示完成跨任务推理的零样本和少样本能力[4]。

在此基础上，Acosta 等于 2022 年提出多模态生物医学 AI（multimodal biomedical AI）概念框架，把精准组学健康、数字化临床试验、远程监护、大流行病监测、数字孪生与虚拟健康助手六类应用纳入统一视角[5]；Moor 等于 2023 年提出通用医疗 AI（generalist medical AI, GMAI）框架，将通用医疗模型的三大核心能力概括为动态任务指定、灵活多模态输入输出与医学领域知识表示[6]。这两套框架共同确立了阶段二的目标函数，不再局限于在单一任务上达到医生水平，而是追求以同一组权重在更广泛任务谱系上保持可用性。

若干代表性体系也在这一阶段集中出现。语言侧，GatorTron 以 820 亿临床笔记词、60 亿 PubMed 词和 25 亿 Wiki 词从头训练至 89B 参数，在临床推理任务上比 BioBERT 提升约 9.5--9.6\%[9]；BioGPT 在 1,500 万 PubMed 摘要上从头训练 GPT-2 Medium，PubMedQA 准确率达到 78.2\%[10]。多模态侧，LLaVA-Med 通过 600K 图文对齐加 60K 自指令对话的两阶段课程，仅使用 8 张 A100、在不足 15 小时内完成训练，在生物医学视觉问答任务上达到与远大于自身的模型相当的性能[11]；Med-Flamingo 以 OpenFlamingo + LLaMA-7B 为骨架，构造从 4,721 本医学教科书抽取的 MTB 数据集（80 万图、5.84 亿 tokens），首次在医学多模态领域验证了少样本上下文学习的可行性[12]。

\subsection{阶段三：生成式 AI 与智能体（2023--2026）}

2022 年末 ChatGPT 引发广泛关注后，医疗 AI 的方法学焦点迅速由基础模型本身转向生成式接口、工具调用与推理增强的组合。Med-PaLM 系列是这一阶段的代表性起点。Singhal 等在 540B Flan-PaLM 上仅通过指令提示微调，更新约 65 个软提示参数，就把科学共识一致率从 61.9\% 提升到 92.6\%，接近临床医生的 92.9\%[8]；Med-PaLM 2 在 PaLM 2 基础上引入集成精炼（ensemble refinement）与检索链（chain of retrieval）两项推理增强机制，使 MedQA 准确率达到 86.5\%，并在 9 个临床轴中的 8 个上获得医生评审者偏好[13]。

Med-Gemini 进一步拓展了基于不确定性的自主搜索能力。当模型输出 token 的 Shannon 熵超过阈值时，系统自动触发网络检索，把外部知识接地纳入推理回路，同时利用 Gemini 1.5 的 1M+ token 上下文窗口，使模型能够直接解析完整电子病历或手术视频，无需依赖 RAG 切片[7,16]。AMIE 则把自博弈与模拟学习作为创新点：在 PaLM 2 基础上，内循环通过 in-context 批评反馈优化对话，外循环通过迭代微调更新权重，模拟覆盖 5,230 种疾病的多轮问诊场景，最终在 159 例 OSCE 风格双盲交叉研究中以 30/32 个专家维度与 25/26 个患者维度显著优于 20 名全科医生[17]。

该阶段的技术路径呈相互促进的发展态势。一方面，推理增强开始进入医疗领域，与传统的指令微调形成互补；另一方面，工具调用与智能体协议让模型从纯权重模型转向带工具的智能体。O'Sullivan 等于 2026 年发表的心脏专科 AMIE，以 Gemini 2.0 Flash 为底座，配套领域工具与多步推理，仅基于 9 例真实病例进行开发，即完成了 107 例 HCM 与心肌病的随机交叉部署，显著错误率较单独专家组下降近一半[19]。Teo 等在 2025 年的综述中将这一阶段的方法学谱系概括为：合成数据系统、扩散模型、规则型 AI、LLM、多模态基础模型、推理与智能体，以及模型蒸馏[7]。

\subsection{技术里程碑时间表}

表 \ref{tab:milestone} 给出 2017 年至 2026 年医疗 AI 关键里程碑的时间表。纵向阅读对应三阶段演进，横向阅读对应训练范式的多样化。

\begin{table}[htbp]
\centering
\caption{医疗 AI 关键里程碑（2017--2026）}
\label{tab:milestone}
\scriptsize
\begin{tabularx}{\textwidth}{L{1.6cm} L{4.5cm} L{2.3cm} Y}
\toprule
时间 & 里程碑 & 所属阶段 & 主要贡献 / 意义 \\
\midrule
2017 & Transformer[A] & 阶段二基础设施 & 跨模态统一架构基础 \\
2018--2019 & BERT / BioBERT / ClinicalBERT[A] & 阶段一$\rightarrow$二 & 生物医学双向编码器，迁移学习起点 \\
2020 & GPT-3 / CLIP / AlphaFold[4,A] & 阶段二基础设施 & LLM、视觉对比学习与蛋白结构三大基石 \\
2021 & Med-BERT[30] / CEHR-BERT[31] & 阶段一$\rightarrow$二 & 结构化 EHR 预训练，奠定专科基础模型路线 \\
2022 & GatorTron[9] / BioGPT[10] / Multimodal Biomedical AI[5] & 阶段二 & 域内规模定律首次验证，多模态综述框架成形 \\
2022 & ChatGPT[A] & 阶段二$\rightarrow$三 & 引爆生成式医疗 AI 范式 \\
2023 & GMAI[6] / Med-PaLM[8] / LLaVA-Med[11] / Med-Flamingo[12] / RETFound[32] & 阶段三起点 & 通用医疗 AI 框架、医学 LLM 与多模态 VLM 并行突破；首个眼科自监督基础模型 \\
2023 & HuatuoGPT[33] / BianQue[34] / DISC-MedLLM[35] / SAM-Med2D[36] & 阶段三 & 中文医疗 LLM 路线启航；分割基础模型扩展至 10 模态 \\
2024 & Med-PaLM 2[13] / Med-Gemini[7,16] / BiomedGPT[15] / MAIRA-2[24] / CONCH[37] / CT-FM (Pai)[38] / MedSAM[39] / EquityMedQA[40] & 阶段三 & 多模态、长上下文、不确定性引导并进；专科基础模型扩展至病理与 CT；公平性工具箱出现 \\
2024 & Foresight 2[41] / AMIE[17] & 阶段三 & EHR 时序专科 LLM；自博弈对话诊断 \\
2025 & MetaMedQA[2] / MedHELM[18] / FUTURE-AI[42] & 阶段三评估转向 & 元认知评估、整体性评估与国际可信 AI 共识 \\
2026 & 心脏专科 AMIE[19] / PreA RCT[20] / 心理治疗认知层[21] / Bean RCT[26] & 阶段三临床证据 & 真实临床随机对照试验证据集中涌现 \\
\bottomrule
\end{tabularx}
\end{table}

\noindent\textit{A：} 时间线背景节点引自 Acosta 等[5]、Moor 等[6]、Thirunavukarasu 等[22]、Teo 等[7]的综述以及通用 AI 文献基础。

\subsection{章节小结}

医疗人工智能经历了多代方法学的演进和叠加。阶段一的对比学习经 CLIP / ConVIRT 演化，成为阶段二多模态对齐的主流组件；阶段二的 BERT 风格双向编码器在阶段三中以检索召回端或专科编码器的形式继续发挥作用，Med-BERT、CEHR-BERT、RETFound、CONCH 等专科基础模型正是这一脉络在 2021--2024 年的延续；阶段一确立的回顾性验证范式则至今仍作为绝大多数医疗 AI 论文默认的评测协议在使用。换言之，医疗 AI 的范式更替并不是简单替代，而是前代技术在后续框架中的重组和再激活。

\section{医疗大语言模型的能力图景}

医疗大语言模型在过去三年完成了一次关键过渡：从“通用底座加领域微调”的基础组合，演进为“通用底座、推理时计算与工具调用”三者的结合体。主流体系分化为指令微调与提示策略、不确定性引导与推理增强、自博弈与模拟学习三条相互渗透的路线。英文主导的体系（Med-PaLM 家族、Med-Gemini、AMIE）在 USMLE 风格的静态基准上已接近性能上限；中文主导的体系（HuatuoGPT、BianQue、PULSE、DISC-MedLLM）则把完成一次合规的多轮问诊作为产品形态，在主动追问、共情表达与本地诊疗指南覆盖上发展出不同于英文路线的方法学侧重。

\subsection{路线一：指令微调与提示策略}

这一路线的核心理念是通过医学领域对齐，最大化通用基础模型已有语言能力的收益。其代表性序列由 Med-PaLM[8] 与 Med-PaLM 2[13] 构成。Singhal 等在 540B Flan-PaLM 之上仅通过指令提示微调，更新约 65 个软提示参数，就把模型在 MedQA（USMLE）上的准确率由 50.3\% 提升到 67.6\%，并把科学共识一致率由 61.9\% 推高到 92.6\%，接近临床医生的 92.9\%[8]。同期由 Med-PaLM 团队构建的 \textbf{MultiMedQA} 评估集整合 MedQA、MedMCQA、PubMedQA、MMLU 临床子集、LiveQA、MedicationQA 与 HealthSearchQA 共计 3,173 题，成为这一阶段事实上的标准基准；同时引入科学共识一致性、不当内容、遗漏、危害程度与偏见五个维度的临床医生人工评估框架。

Med-PaLM 2 在 PaLM 2 基础上叠加两项关键推理增强。\textbf{集成精炼} 通过多路并行推理后综合输出，降低单一采样的偶然性；\textbf{检索链} 则把初始答案拆分为若干可验证子声明，逐项调用搜索完成事实接地。这两项机制使其在 MedQA 上达到 86.5\% 的准确率，并在 9 个临床轴中的 8 个上获得医生评审者偏好[13]。

\subsection{路线二：不确定性引导与推理增强}

当指令微调把静态准确率推近基准上限后，模型的剩余错误主要表现为 \textbf{过度自信的事实捏造}。Med-Gemini 系列以此为切入点，将训练目标从单纯提升正确率转向“在不确定时主动检索”[7]。其核心机制是基于输出 token 的 Shannon 熵触发外部搜索：当模型对下一 token 的预测分布过于平坦时，自动调用网络检索，把外部知识接地纳入推理回路。Med-Gemini 在 MedQA（USMLE）上取得 91.1\% 的准确率；在团队对评测题重新标注、修复错答与歧义后，SOTA 准确率上移至 91.8--92.9\%，这表明 \textbf{评测瓶颈已在一定程度上由模型能力转向标注质量}，是阶段三评估范式转向的早期信号。

Med-Gemini 的另一项关键设计是利用 Gemini 1.5 的 1M+ token 上下文窗口直接解析完整电子病历或手术视频，将检索增强生成中的切片、召回与拼接步骤大幅简化为长上下文输入[7]。这一设计在概念上把“模型加 RAG”的两段式架构压缩为“模型、上下文与工具调用”的三层架构，是从语言驱动过渡到工具与反馈驱动的具体落地范例。其 2D / 3D / Polygenic 扩展[16] 则展示了这一架构在多模态长尾任务上的延展能力。

\subsection{路线三：自博弈与模拟学习}

第三条路线以 AMIE 为代表[17]。其方法学创新在于用自博弈绕开真实医患对话数据在隐私与规模上的约束：在 PaLM 2 基础上，模型同时扮演医生与模拟患者两个角色，内循环通过 in-context 批评机制对每一轮对话即时反馈，外循环则把累积的高质量对话作为新的微调数据迭代更新权重，模拟覆盖 5,230 种疾病的多轮问诊场景。在 159 例 OSCE 风格双盲交叉研究中，AMIE 在 30/32 个专家维度与 25/26 个患者维度上显著优于 20 名全科医生；在鉴别诊断推理和管理建议合理性两个维度上，其优势超过 0.5 个李克特量级。

AMIE 将 \textbf{模拟学习的合规边界} 纳入主流方法学讨论。与依赖真实医患对话训练的中文路线相比，自博弈能在不接触患者隐私的前提下生成功能近似的训练信号；但其代价是模拟分布与真实分布之间潜在的偏差。

\subsection{中文医疗大语言模型：另一种问题设定}

英文体系通常将“通过 USMLE”或“增益医生工作流”作为里程碑，因此更强调应答能力；中文体系则将“完成一次合规的多轮问诊”作为产品形态，因此更强调主动追问能力。这一差异催生了方法学层面的独立创新。

HuatuoGPT 是该路线的代表性早期工作[33]。Zhang 等观察到，当时医疗 LLM 只有“纯粹蒸馏自 ChatGPT”与“只在真实医患对话上训练”两种路径，前者缺乏专业诊断能力，后者则回复简短、不够友好。他们提出混合数据 SFT 加混合反馈强化学习（RLMF）的双重融合架构：前者同时利用 ChatGPT 蒸馏数据与真实医患对话数据，后者由 ChatGPT 与人类医生联合提供奖励信号。HuatuoGPT 基于 LLaMA-13B（Ziya-LLaMA 初始化）训练，在多轮对话评估中相对 DoctorGLM 胜率为 97\%，相对 ChatGPT 胜率为 62\%，并在 cMedQA2、webMedQA、Huatuo-26M 等中文医疗 QA 基准上取得领先表现。

BianQue 进一步强调模型主动追问的能力，提出 \textbf{Chain of Questioning（CoQ）} 概念[34]。Chen 等指出，真实临床场景中医生需通过多轮问诊逐步了解患者状况，才能给出个性化建议；而当时的中文医疗 LLM 因训练数据缺乏多轮追问样本而“只会建议、不会提问”。BianQue 团队构建了 243 万条多轮健康对话数据集 BianQueCorpus，通过 ChatGPT 对真实医患对话中医生的建议进行润色，保持问题与建议的比例约为 46.2\% 对 53.8\%。以 ChatGLM-6B 为底座做全参数微调后，BianQue 在 MedDG 数据集上 BLEU-1 达到 14.86、ROUGE-L 达到 19.56，自定义的 PQA 指标达到 0.81，均显著高于 ChatGPT 基线。

PULSE 与 DISC-MedLLM 进一步扩展了中文医疗人工智能的规模。PULSE 由上海人工智能实验室 OpenMEDLab 团队发布，以 BLOOMZ-7B 与自研中文基座扩展版本为底座，在约 400 万条中文医学指令数据上做监督微调，在 MedQA-MCMLE 与 CMB 上压制同期开源中文通用 LLM[43]。DISC-MedLLM 则由复旦大学团队基于 Baichuan-13B-Base 构建，DISC-Med-SFT 数据集通过三条管线生成约 47 万条样本，其多轮一致性与行为得分超过 HuatuoGPT、BianQue 与中文 GPT-3.5[35]。

中文路线的评估生态也在同步成熟。Wang 等于 NAACL 2024 主会发表的 \textbf{CMB} 覆盖六大临床工种、中医、药剂、护理等子领域，并按中国医师资格考试结构设计层级化任务[44]；其后的 CMExam 与 MedBench 进一步加入了临床安全维度。CMB 与 MedBench 中包含中医与医保政策等本地化任务，使中文医疗基准与 MultiMedQA、MedHELM 形成了结构性差异。

\subsection{章节小结}

表 \ref{tab:llm-compare} 将上述四条路线的代表性模型并置对比，可以看到同一阶段内的方法学多样性远大于阶段之间的差异。这表明 \textbf{当下医疗 LLM 的能力上限不再由底座大小单独决定，而是由训练目标、对齐数据与推理时计算三个因素共同决定}。

\begin{table}[htbp]
\centering
\caption{代表性医疗大语言模型横向对比}
\label{tab:llm-compare}
\scriptsize
\begin{tabularx}{\textwidth}{L{1.7cm} L{0.9cm} L{2cm} L{2.6cm} Y Y L{1.8cm}}
\toprule
模型 & 年份 & 主干 & 训练 / 对齐数据 & 关键贡献 & 代表性结果 & 目标场景 \\
\midrule
BioGPT[10] & 2022 & GPT-2 Medium 347M & 1,500 万 PubMed 摘要 & 首个生物医学生成式预训练模型 & PubMedQA 78.2\% & 文献问答 \\
GatorTron[9] & 2022 & Megatron-LM BERT & 820 亿临床笔记词 & 89B 临床域内全参预训练 & NLI 90.20\%，MQA 93.10\% & 临床 NLP \\
Med-PaLM[8] & 2023 & Flan-PaLM 540B & 指令提示微调（约 65 软提示） & 多维医生评估与 MultiMedQA & MedQA 67.6\%，危害比例 5.9\% 近似医生 5.7\% & USMLE 风格问答 \\
Med-PaLM 2[13] & 2025 & PaLM 2 & MedQA / MedMCQA / HSQA & 集成精炼与检索链 & MedQA 86.5\%，9 轴中 8 轴胜医生 & 工作流问答 \\
Med-Gemini[7] & 2024 & Gemini 1.0 Ultra & 多模态混合与自训练 & 不确定性引导搜索与 1M token 上下文 & MedQA 91.1\%，重标后 91.8--92.9\% & 长上下文与工具调用 \\
AMIE[17] & 2025 & PaLM 2 & 真实对话与自博弈 & In-context 批评与迭代微调 & OSCE 30/32 专家维度优于 GP & 多轮诊断对话 \\
HuatuoGPT[33] & 2023 & LLaMA-13B（Ziya） & Huatuo-26M + Med-Dialog + ChatGPT 蒸馏 & 混合数据 SFT 与 RLMF & 相对 DoctorGLM 胜率 97\%，相对 ChatGPT 62\% & 中文多轮咨询 \\
BianQue[34] & 2023 & ChatGLM-6B & BianQueCorpus 243 万 & Chain of Questioning & MedDG PQA 0.81，对比 ChatGPT 0.63 & 主动追问问诊 \\
PULSE[43] & 2023 & BLOOMZ-7B + 自研 20B & 约 400 万中文医学指令 & 中文医考、中医与指南覆盖 & CMB / MedQA-MCMLE 优于同期开源中文通用 LLM & 中文医考与咨询 \\
DISC-MedLLM[35] & 2023 & Baichuan-13B & DISC-Med-SFT 47 万 & 三管线数据与偏好对齐 & 多轮一致性优于 HuatuoGPT、BianQue 与中文 GPT-3.5 & 中文医疗对话 \\
\bottomrule
\end{tabularx}
\end{table}

\section{多模态医疗大模型}

多模态医疗大模型的演进围绕“如何把视觉、生物信号、组学等模态稳健接入语言模型，并在下游任务上保持统一权重的可用性”展开。目前三条主流融合路线在 MultiMedBench 等综合基准上的性能差距已小于 5 个百分点，竞争焦点也随之转移至三类不对称，即模态可得性不对称、评估指标不对称与空间接地不对称。

\subsection{架构路线一：联合融合}

联合融合将视觉编码器输出的 token 与文本 token 直接拼接，再由单一解码器在交错序列上完成跨模态推理。Med-PaLM M 是这条路线在医疗领域最具代表性的落地[14]：其以 PaLM-E 为骨架，把 12B、84B、562B 三档参数规模分别在 MultiMedBench 的 14 个任务上做统一权重训练，覆盖 MedQA、MedMCQA、PubMedQA、MIMIC-III RRS、VQA-RAD、Slake-VQA、Path-VQA、MIMIC-CXR、PAD-UFES-20、VinDr-Mammo、CBIS-DDSM、PrecisionFDA V2 等基准。其 562B 版本在 14 个任务中全部达到或接近 SOTA，是当时通用医疗多模态模型的能力上限。Med-PaLM M 提供的超过 100 万样本的多模态医疗任务集合，也长期构成这一方向覆盖最广的统一评估基准。

\subsection{架构路线二：门控交叉注意力}

门控交叉注意力通过 Perceiver Resampler 把视觉特征压缩到固定槽位，再以门控交叉注意力将视觉信息注入 LLM 中间层，在引入视觉模态的同时尽量保留主干语言模型的原有能力。Med-Flamingo 将这一思路引入医学领域[12]，以 OpenFlamingo + LLaMA-7B 为骨架（共 8.3B 参数），从 4,721 本医学教科书抽取并构建出 MTB 数据集（80 万图、5.84 亿 tokens）用于继续训练。该模型首次在医学多模态领域验证了 \textbf{少样本上下文学习} 的可行性，使模型能够在 4--8 个标注样本的提示下完成 PathVQA、VQA-RAD 等任务的零样本到少样本推理。

\subsection{架构路线三：离散 token 化与统一序列建模}

离散 token 化将图像、文本、表格乃至空间位置统一映射为离散 token，再由 Transformer 在共享词表上完成序列建模，把跨模态统一性推进到更彻底的形式。BiomedGPT 是这一路线的轻量化代表[15]：其基于 OFA 框架，用 VQ-GAN 把图像离散为视觉 token，与文本 BPE token 共享一个 59,457 维统一词表；33M、93M、182M 三档参数规模在 25 项任务中有 16 项达到 SOTA。尤其值得强调的是，182M 参数版本相对 562B 的 Med-PaLM M 在参数规模上缩小超过 3,000 倍，却仍在多项放射、病理与皮肤镜任务上保持接近 SOTA 的性能，说明模型规模并非能力上限的唯一决定因素。

MAIRA-2 把离散 token 化进一步扩展到 \textbf{空间坐标} 维度[24]。其以 100$\times$100 网格坐标 token 编码病灶空间定位，使模型能够生成带句级边界框的接地放射报告。MAIRA-2 由 Vicuna 与 Rad-DINO 组合构建（约 7B 参数），在 MIMIC-CXR 上取得 BLEU-4 23.1、RadGraph-F1 34.6 等 SOTA，专家评审显示 91\% 的句子可以被临床医生直接接受。

\subsection{长上下文与模态扩展}

除上述三条主流架构外，Med-Gemini 系列还以 \textbf{长上下文窗口} 为多模态扩展提供了另一条路径[7,16]。Med-Gemini-M 1.5 支持 1M+ token 上下文，这意味着模型可以将完整电子健康记录、长视频或多张影像直接纳入提示上下文，而不再依赖 RAG 切片。Med-Gemini-2D 把 X 光、皮肤镜、眼底等 2D 模态接入；Med-Gemini-3D 则把头颈 CT、肺部 CT 等三维体素数据通过切片采样接入并完成报告生成；Med-Gemini-Polygenic 进一步把基因组多基因风险评分以零样本方式纳入预测。需要指出的是，Med-Gemini-3D 在结构化报告生成任务上仍存在显著幻觉，而 CT-US1 等私有数据集的存在也使其结论的可复核性受限。

\subsection{接地与可追溯：以 RadFact 为例}

多模态医疗报告生成长期以 BLEU、ROUGE、CIDEr 等基于词汇重叠的指标作为评估标准，但这些指标对医学报告中的临床显著性错误并不敏感。MAIRA-2 配套提出的 RadFact[24] 是针对这一缺陷的系统响应。该框架利用 Llama-3-70B 对模型生成报告与参考报告做句级逻辑蕴含推断，并给出四维评估：逻辑精确率与召回率，以及空间精确率与召回率。在 MIMIC-CXR 上，MAIRA-2 的 RadFact 指标比 BLEU 等传统指标具有更强的区分度，并与放射科专家评分高度相关。这表明多模态医疗模型的评估正在从“词汇重叠”转向“临床接地与可验证性”。

\subsection{章节小结}

多模态医疗大模型的瓶颈已从“如何把视觉接入 LLM”转向“如何让有限的数据、评估与接地基础设施跟上模型能力”。具体而言，这一瓶颈在三个维度上呈现不对称：第一，CXR、皮肤镜、眼底等 2D 影像在公开数据中占主导，但病理、3D CT、心脏超声、内镜视频等模态仍处于长尾；第二，BLEU 与 ROUGE 仍是大多数论文的默认指标，而 RadFact 这类临床显著性加权评估只在少数体系中实现协议化；第三，句级边界框接地已在胸片报告中达到较高成熟度，但其他模态和任务中同等级别的接地能力仍属空白。

\begin{table}[htbp]
\centering
\caption{代表性多模态医疗大模型横向对比}
\label{tab:mm-compare}
\scriptsize
\begin{tabularx}{\textwidth}{L{2.1cm} L{2.8cm} L{1.4cm} L{2.2cm} Y Y}
\toprule
模型 & 模态覆盖 & 规模 & 架构路线 & 突出贡献 & 关键局限 \\
\midrule
LLaVA-Med[11] & 文本 + 影像（CXR / CT / MRI / 病理） & 7B / 13B & 联合融合（视觉投影） & GPT-4 自指令与两阶段课程；8$\times$A100 训练少于 15 小时 & 输入仅 224$\times$224，幻觉显著 \\
Med-Flamingo[12] & 文本 + 影像 & 8.3B & 门控交叉注意力 & MTB 80 万图，验证少样本 ICL & PathVQA 较弱，仍属 PoC \\
Med-PaLM M[14] & 文本 + CXR + 病理 + 皮肤 + 钼靶 + 基因组 & 12B / 84B / 562B & 联合融合（PaLM-E） & MultiMedBench 14 任务统一权重 & 未开源，缺前瞻验证 \\
BiomedGPT[15] & 文本 + 2D 图像 + 表格 & 33M / 93M / 182M & 离散 token 化 & 全开源且轻量化，相对 562B 缩小 3,000 倍以上 & 零样本 VQA 低于 60\% \\
Med-Gemini-2D/3D/Polygenic[16] & 2D 影像 + 3D CT + 基因组 & 未披露 & 长上下文 + 多模态 & 3D 报告生成与 PRS 零样本预测 & 3D 幻觉明显，CT-US1 私有 \\
MAIRA-2[24] & CXR + 历史片 + 临床文本 & 约 7B & 离散 token（空间坐标） & 句级边界框接地与 RadFact 评估 & 仅覆盖 CXR，51\% 训练数据私有 \\
\bottomrule
\end{tabularx}
\end{table}

\section{专科基础模型的不可替代位置}

虽然通用基础模型不断刷新综合基准，但在结构化 EHR 时序、专科影像、组学预测等具备强结构先验的子领域，专科自监督基础模型在精度、样本效率与跨机构泛化三个维度上系统性优于通用多模态大模型。通用与专科并非替代关系，而更接近互补关系。

\subsection{结构化 EHR 时序：从 Med-BERT 到 Foresight 2}

结构化 EHR 是医疗人工智能中最具结构先验的子领域。它本质上是一组以患者为单位、按时间排序的 ICD、SNOMED、CPT 与药物编码序列，与自然语言相比具有三个关键结构特征：\textbf{编码空间离散且封闭、时间间隔关键、患者间方差远大于编码方差}。这些特征决定了直接套用通用 LLM 在 EHR 上很难取得最优表现。

Med-BERT 与 CEHR-BERT 分别代表结构先验工程化的两条早期路径[30,31]。Med-BERT 在 Cerner Health Facts 的 2,849 万患者诊断码序列上预训练 17M 参数的 BERT 风格模型，关键创新是在 masked LM 之外引入延长住院时长预测作为领域特定的第二预训练目标；在 300 个样本的小数据场景下，Med-BERT 接 Bi-GRU 仍可达 AUC 0.75，相当于训练集扩大约 10 倍后的 Bi-GRU 性能。CEHR-BERT 则聚焦 \textbf{时间结构}，以人工时间标记把相邻就诊间隔以离散 token 形式插入序列，配合概念、时间与年龄嵌入，并以就诊类型预测替换 BERT 原版 NSP；在仅用 5\% 训练数据微调的情况下，CEHR-BERT 仍能超越使用全量数据训练的对比模型。这两项工作共同揭示出：\textbf{样本效率会随结构先验显著放大}。

Foresight 2 则进一步把结构先验与通用语言能力结合起来[41]。其关键设计是 \textbf{语境化患者时间线}：保留每个 SNOMED 概念被提及的原始文本上下文，并把 SNOMED 编码嵌入 LLM 分词器，让 7B 规模的通用 LLM（Mistral-7B 与 LLaMAv2-7B）能够高效处理标准化医学本体的输入与输出。在 MIMIC-III 535 名患者的风险预测测试集上，Foresight 2 的 P@5 达到 0.90，显著高于 GPT-4-turbo 的 0.65。这一对比直接说明：当 EHR 任务被恰当地结构化建模时，\textbf{7B 专科模型可以系统性优于参数规模更大的通用模型}。同样重要的是，Foresight 2 的消融实验显示，移除患者时间线中的文本上下文后性能下降约 40\%，说明结构先验与通用语言能力必须同时具备。

\subsection{影像与分割专科基础模型}

专科影像领域的方法学创新沿三条线索展开：\textbf{视觉自监督表示学习、视觉语言对齐、参数高效适配}。这三条线索分别对应不同的下游任务结构与数据可得性约束，但都在小样本与长尾场景下展现出对通用模型的系统性优势。

视觉自监督一脉以 RETFound 为旗舰[32]。其以 ViT-Large 在 90.4 万张彩色眼底照片与 73.6 万张 OCT 图像上做 MAE 风格预训练，下游覆盖糖网分级、青光眼、AMD 等 8 项眼科诊断任务以及心衰、心梗、缺血性卒中、帕金森病等 4 项三年预后预测；全部 12 项任务的 AUROC 均优于 ImageNet 预训练与监督基线，且 \textbf{优势在每类样本少于 100 时最为显著}。Pai 等进一步把同一思路扩展到 3D CT，在跨机构外部验证集上 AUC 与 C-index 显著优于 ImageNet-3D 与从头训练基线[38]。

视觉语言对齐一脉的代表是 CONCH[37]。其方法学包含两步：先以 1,600 万张病理 ROI 做 iBOT 风格视觉自监督训练，再在 117 万对病理图像--文本上以 CoCa 框架做图文对比与描述生成联合训练。这种 \textbf{视觉表示与语言接地分层使用} 的设计让 CONCH 在 14 类基准上有 11 项达到 SOTA，且在 zero-shot 罕见癌种分类上较 PLIP 提升超过 10 个百分点。

参数高效适配一脉则以 SAM-Med2D 与 MedSAM 为代表[36,39]。二者均以 Meta SAM 为骨架，通过 adapter、continual pretraining 等手段做专科适配。SAM-Med2D 合并 31 个公开数据集，在 9 类外部数据集上 Dice 较原生 SAM 提升 10--30 个百分点；MedSAM 则在不同数据规模与训练协议下复现了相同思路。这条路线以通用基础模型为骨架、以参数高效手段做专科适配，在精度与经济性之间取得了现实平衡。

\subsection{通用模型局限性的对照证据}

Jin 等发表在 \textit{npj Digital Medicine} 的研究为通用多模态模型在同类任务上的局限提供了直接对照证据[45]。他们以多类医学影像任务评估 GPT-4V 的多模态推理能力，发现一个反复出现的失效模式：\textbf{模型的选项准确率并不低，但图像感知与推理路径中包含大量隐性错误}。也就是说，模型往往凭借选项线索给出正确答案，但其对图像的实际描述与图中内容不符，或者关键区域识别错误。由此可见，单看选择题准确率会高估通用多模态模型在临床部署中的可靠性。

\subsection{章节小结与发展图景}

Foresight 2 在 EHR 上压制 GPT-4-turbo，RETFound 在小样本眼科任务上优于 ImageNet-ViT，CONCH 在 zero-shot 罕见癌种任务上优于 PLIP，Pai 等在小样本预后任务上优于 ImageNet-3D。综上，在结构化、长尾、小样本三类特征同时出现的子领域，\textbf{专科自监督基础模型仍然不可替代}。

综合前文可见，能力扩张的来源是数据规模、训练目标与结构先验三者的乘积，没有任何一项能够单独决定上限。当下“通用多模态大模型 + 专科基础模型 + 检索与工具调用”的混合架构，是这一时期最具代表性的事实图景。

\begin{table}[htbp]
\centering
\caption{通用模型与专科基础模型在四类任务上的相对优势矩阵}
\label{tab:specialist-matrix}
\scriptsize
\begin{tabularx}{\textwidth}{L{2.3cm} L{2.5cm} Y Y L{2cm} L{1.6cm}}
\toprule
任务类别 & 代表性任务 & 通用大模型 & 专科基础模型 & 相对优势归属 & 证据来源 \\
\midrule
结构化时序预测 & EHR 风险预测、再入院、HF 预后 & GPT-4-turbo P@5 0.65 & Foresight 2 P@5 0.90；Med-BERT 小样本提升约 20\% AUC & 专科压制 & [41,30,31] \\
罕见类别长尾分类 & 罕见癌种 zero-shot、罕见病眼科 & PLIP、CLIP zero-shot 较弱 & CONCH zero-shot 提升超过 10 个百分点；RETFound 小样本 SOTA & 专科压制 & [37,32] \\
高分辨率影像几何 & 3D CT 报告、病理 WSI、跨模态分割 & GPT-4V 存在隐性感知失效；Med-Gemini-3D 幻觉明显 & CONCH、Pai 等、MedSAM 跨机构更稳健 & 专科压制 + 通用辅助 & [38,39,45] \\
开放语义多轮交互 & 多轮诊断对话、文献综述、患者教育 & AMIE / Med-PaLM 2 在 OSCE 等任务中占优 & 专科 LLM 缺失或较弱 & 通用压制 & [17,13] \\
\bottomrule
\end{tabularx}
\end{table}

\section{讨论}

前文描绘了一幅以能力扩张为主线的发展图景，但除能力本身之外，临床证据、评估方法与治理框架三方面仍呈现出明显缺口。

第一类问题集中在人机协作的非线性。Goh 等招募 50 名内科医师在 6 个真实病例上做诊断推理，GPT-4 独立组评分 92\%，但“医生 + GPT-4”组仅 76\%，与传统资源组无显著差异[25]；Bean 等在 1,298 名公众用户的预注册随机对照试验中观察到更剧烈的剪刀差：三款主流 LLM 在脱离用户的封闭设置下识别 10 类病症的准确率为 90.8--99.2\%，但真实公众用户经 LLM 辅助后的病症识别率反而显著低于互联网搜索组[26]。反向证据则来自约束化部署：Pais 等的 MEDIC 将 LLM 严格限定为结构化抽取器，在在线药房前瞻部署中使近失误事件下降 33\%[27]；PreA 在两家三甲医院的随机对照试验中把问诊时长缩短 28.7\%、跨科室协调度提升 113.1\%[20]；心脏专科 AMIE 将显著错误率从 24.3\% 降到 13.1\%[19]。这些证据共同说明，模型独立基准成绩与真实临床效用之间存在系统性的“协作折扣”，而工作流约束、用户素养与共同设计强度，比模型规模本身更能解释这一折扣的方向与大小。

第二类问题在于评估方法学的拓扑滞后。MedHELM 与 29 名临床医生联合构建了 5 大类、22 子类、121 个临床任务及配套 37 个基准，发现所有模型在行政与工作流类别得分最差，恰好对应第一类问题中最易出问题的领域[18]。MetaMedQA 把“我不知道”作为显式选项加入 MedQA 扩展集，结果 12 个主流模型中有 9 个 Unknown Recall 为 0[2]；Kim 等构造的对抗推理题表明，人类医生平均 66\%，最优 LLM 仅 52\%，而医疗专用模型在该任务上甚至接近 0\%[23]。EquityMedQA 与 RadFact 则分别在公平性与多模态接地两个维度上引入新的评估面[24,40]。合并这些结果可以看到，评估科学的当务之急不在于构造更难的题目，而在于扩展评估空间的拓扑，包括任务覆盖、不确定性、公平性、空间接地、被试身份与报告规范等多个维度。

第三类问题潜伏于训练--治理链路的隐性脆弱。Alber 等的数据投毒实验表明，只需在通用 LLM 预训练语料中替换 0.001\% 的训练 token，便足以让 4B 模型在医疗有害内容生成率上显著上升，而所有被投毒模型在 5 项主流医疗基准上的得分与未投毒基线无显著差异[29]。这说明，任何依赖标准基准做上市前评估的监管体系，都可能对这类攻击结构性失效。该脆弱性又与训练数据透明度不足相互放大：系统综述显示 87.7\% 的医疗 LLM 为闭源体系[1]，而 Med-PaLM M、Med-Gemini、MAIRA-2 等代表性模型的训练数据多数私有。监管框架的滞后则被多篇综述和共识反复指出：现行 FDA 与 NMPA 监管基于固定权重，与持续学习、智能体调用、自我更新的 LLM 范式不匹配[3,7,42]。因此，训练数据透明度应当被视为医疗 AI 安全性的必要条件，而不是可选的良好实践。

总体而言，医疗 AI 的能力扩张本身不足以自动保证临床效用、评估有效性与治理合规。患者直面部署还需要超出基准评估的额外维度，例如可读性、健康素养与文化适配[28]。下一阶段的进步很可能不再主要来自更大的模型，而来自把模型与专科底座、外部验证器以及持续审计层组装为可工程化、可监管化的系统。

\section{结论}

医疗人工智能在过去四年经历了基础模型化、生成式化与多模态化三个层面的发展。Med-PaLM 2、Med-Gemini、AMIE 等通用与多模态大模型在多个静态基准上已超过医生水平，并在 OSCE、心脏专科、心理治疗、转诊等真实场景中取得了首批随机对照试验证据；LLaVA-Med、Med-Flamingo、BiomedGPT 等模型共同绘制出覆盖语言、影像、组学的多模态能力图景；Foresight 2、RETFound、MedSAM 等专科基础模型则在 EHR 时序、眼底、病理、CT、分割等子领域上证明结构先验不可替代。

这一演进并非范式间的相互替代，而是层级化的叠加复用。对比学习与双向编码器延续进入生成式阶段，基础模型骨架与多模态对齐组件又以专科底座与工具调用的形式在第三阶段被再度激活。能力扩张同样并非单一变量驱动：训练目标的创新可以在不增加底座规模的前提下解锁新能力，结构先验在小样本与长尾任务上具有数量级意义的优势，而通用化与专科化是必须同时推进的两个方向。

然而，能力曲线的持续上扬并未带动三条关键侧线的同步成熟：临床协作的非线性效应、评估方法学的拓扑滞后，以及训练--治理链路的隐性脆弱，共同构成下一阶段亟待补齐的结构性缺口。医疗人工智能的发展图景已经初具规模，但距离“可用、可信、可监管”的目标仍有相当距离。缩小这一距离的关键，不再仅取决于模型本身，而在于围绕模型构建的系统。

\begin{thebibliography}{99}

\bibitem{ref1} Chen, S. F., Oermann, E. K. \textit{et al}. LLM-assisted systematic review of large language models in clinical medicine. \textit{Nature Medicine} \textbf{32}, 1152--1159 (2026). DOI: 10.1038/s41591-026-04229-5.
\bibitem{ref2} Griot, M. \textit{et al}. Large language models lack essential metacognition for reliable medical reasoning. \textit{Nature Communications} (2025). DOI: 10.1038/s41467-024-55628-6.
\bibitem{ref3} Rajpurkar, P., Chen, E., Banerjee, O. \& Topol, E. J. AI in health and medicine. \textit{Nature Medicine} \textbf{28}, 31--38 (2022). DOI: 10.1038/s41591-021-01614-0.
\bibitem{ref4} Bommasani, R. \textit{et al}. On the Opportunities and Risks of Foundation Models. arXiv:2108.07258 (2022).
\bibitem{ref5} Acosta, J. N., Falcone, G. J., Rajpurkar, P. \& Topol, E. J. Multimodal biomedical AI. \textit{Nature Medicine} \textbf{28}, 1773--1784 (2022). DOI: 10.1038/s41591-022-01981-2.
\bibitem{ref6} Moor, M. \textit{et al}. Foundation models for generalist medical artificial intelligence. \textit{Nature} \textbf{616}, 259--265 (2023). DOI: 10.1038/s41586-023-05881-4.
\bibitem{ref7} Teo, Z. L., Thirunavukarasu, A. J. \& Ting, D. S. W. Generative artificial intelligence in medicine. \textit{Nature Medicine} \textbf{31} (2025). DOI: 10.1038/s41591-025-03983-2.
\bibitem{ref8} Singhal, K. \textit{et al}. Large language models encode clinical knowledge. \textit{Nature} \textbf{620}, 172--180 (2023). DOI: 10.1038/s41586-023-06291-2.
\bibitem{ref9} Yang, X. \textit{et al}. A large language model for electronic health records (GatorTron). \textit{npj Digital Medicine} (2022). DOI: 10.1038/s41746-022-00742-2.
\bibitem{ref10} Luo, R. \textit{et al}. BioGPT: generative pre-trained transformer for biomedical text generation and mining. \textit{Briefings in Bioinformatics} \textbf{23}, bbac409 (2022). DOI: 10.1093/bib/bbac409.
\bibitem{ref11} Li, C., Wong, C., Zhang, S. \textit{et al}. LLaVA-Med: Training a Large Language-and-Vision Assistant for Biomedicine in One Day. \textit{NeurIPS 2023 Datasets and Benchmarks}.
\bibitem{ref12} Moor, M., Huang, Q. \textit{et al}. Med-Flamingo: A Multimodal Medical Few-Shot Learner. arXiv:2307.15189 (2023).
\bibitem{ref13} Singhal, K. \textit{et al}. Toward expert-level medical question answering with large language models (Med-PaLM 2). \textit{Nature Medicine} \textbf{31}, 943--950 (2025). DOI: 10.1038/s41591-024-03423-7.
\bibitem{ref14} Tu, T., Azizi, S. \textit{et al}. Towards Generalist Biomedical AI (Med-PaLM M). \textit{NEJM AI} (2024). arXiv:2307.14334.
\bibitem{ref15} Zhang, K. \textit{et al}. A generalist vision-language foundation model for diverse biomedical tasks (BiomedGPT). \textit{Nature Medicine} (2024). DOI: 10.1038/s41591-024-03185-2.
\bibitem{ref16} Yang, L. \textit{et al}. Advancing Multimodal Medical Capabilities of Gemini (Med-Gemini-2D/3D/Polygenic). arXiv:2405.03162 (2024).
\bibitem{ref17} Tu, T., Schaekermann, M., Palepu, A. \textit{et al}. Towards conversational diagnostic artificial intelligence (AMIE). \textit{Nature} \textbf{642} (2025). DOI: 10.1038/s41586-025-08866-7.
\bibitem{ref18} Bedi, S. \textit{et al}. Holistic evaluation of large language models for medical tasks with MedHELM. \textit{Nature Medicine} \textbf{32}, 943--951 (2026). DOI: 10.1038/s41591-025-04151-2.
\bibitem{ref19} O'Sullivan, J. W., Palepu, A., Ashley, E., Tu, T. \textit{et al}. A large language model for complex cardiology care. \textit{Nature Medicine} \textbf{32} (2026). DOI: 10.1038/s41591-025-04190-9.
\bibitem{ref20} Tao, X., Zhou, S., Ding, K. \textit{et al}. An LLM chatbot to facilitate primary-to-specialist care transitions: a randomized controlled trial. \textit{Nature Medicine} \textbf{32}, 934--942 (2026). DOI: 10.1038/s41591-025-04176-7.
\bibitem{ref21} Rollwage, M. \& McFadyen, J. \textit{et al}. A cognitive layer architecture to support large-language model performance in psychotherapy interactions. \textit{Nature Medicine} (2026). DOI: 10.1038/s41591-026-04278-w.
\bibitem{ref22} Thirunavukarasu, A. J. \textit{et al}. Large language models in medicine. \textit{Nature Medicine} \textbf{29}, 1930--1940 (2023). DOI: 10.1038/s41591-023-02448-8.
\bibitem{ref23} Kim, J. \& Bernardo, D. \textit{et al}. Limitations of large language models in clinical problem-solving arising from inflexible reasoning. \textit{Scientific Reports} (2025). DOI: 10.1038/s41598-025-22940-0.
\bibitem{ref24} Bannur, S., Hyland, S. L. \textit{et al}. MAIRA-2: Grounded Radiology Report Generation. arXiv:2406.04449 (2024).
\bibitem{ref25} Goh, E. \textit{et al}. Large language model influence on diagnostic reasoning: a randomized clinical trial. \textit{JAMA Network Open} (2024). DOI: 10.1001/jamanetworkopen.2024.40969.
\bibitem{ref26} Bean, A. M. \& Mahdi, A. \textit{et al}. Reliability of LLMs as medical assistants for the general public: a randomized preregistered study. \textit{Nature Medicine} \textbf{32}, 609--615 (2026). DOI: 10.1038/s41591-025-04074-y.
\bibitem{ref27} Pais, C., Liu, J. \textit{et al}. Large language models for preventing medication direction errors in online pharmacies (MEDIC). \textit{Nature Medicine} (2024). DOI: 10.1038/s41591-024-02933-8.
\bibitem{ref28} Beale, S. K. \textit{et al}. Comparing physician and AI chatbot responses to posthysterectomy questions posted to a public social media forum. \textit{AJOG Global Reports} (2025). DOI: 10.1016/j.xagr.2025.100553.
\bibitem{ref29} Alber, D. A., Oermann, E. K. \textit{et al}. Medical large language models are vulnerable to data-poisoning attacks. \textit{Nature Medicine} (2025). DOI: 10.1038/s41591-024-03445-1.
\bibitem{ref30} Rasmy, L., Xiang, Y., Xie, Z., Tao, C. \& Zhi, D. Med-BERT: pretrained contextualized embeddings on large-scale structured electronic health records for disease prediction. \textit{npj Digital Medicine} \textbf{4}, 86 (2021). DOI: 10.1038/s41746-021-00455-y.
\bibitem{ref31} Pang, C., Jiang, X., Kalluri, K. S. \textit{et al}. CEHR-BERT: Incorporating temporal information from structured EHR data to improve prediction tasks. arXiv:2111.08585 (2021).
\bibitem{ref32} Zhou, Y., Chia, M. A., Wagner, S. K. \textit{et al}. A foundation model for generalizable disease detection from retinal images (RETFound). \textit{Nature} \textbf{622}, 156--163 (2023). DOI: 10.1038/s41586-023-06555-x.
\bibitem{ref33} Zhang, H., Chen, J., Jiang, F. \textit{et al}. HuatuoGPT, Towards Taming Language Models To Be a Doctor. \textit{Findings of EMNLP 2023}. arXiv:2305.15075.
\bibitem{ref34} Chen, Y., Wang, Z., Xing, X. \textit{et al}. BianQue: Balancing the Questioning and Suggestion Ability of Health LLMs with Multi-turn Health Conversations Polished by ChatGPT. arXiv:2310.15896 (2023).
\bibitem{ref35} Bao, Z., Chen, W., Xiao, S. \textit{et al}. DISC-MedLLM: Bridging General Large Language Models and Real-World Medical Consultation. arXiv:2308.14346 (2023).
\bibitem{ref36} Cheng, J., Ye, J., Deng, Z. \textit{et al}. SAM-Med2D. arXiv:2308.16184 (2023).
\bibitem{ref37} Lu, M. Y., Chen, B., Williamson, D. F. K. \textit{et al}. A visual-language foundation model for computational pathology (CONCH). \textit{Nature Medicine} \textbf{30}, 863--874 (2024). DOI: 10.1038/s41591-024-02856-4.
\bibitem{ref38} Pai, S., Bontempi, D., Hadzic, I. \textit{et al}. Foundation model for cancer imaging biomarkers. \textit{Nature Machine Intelligence} \textbf{6}, 354--367 (2024). DOI: 10.1038/s42256-024-00807-9.
\bibitem{ref39} Ma, J., He, Y., Li, F. \textit{et al}. Segment anything in medical images (MedSAM). \textit{Nature Communications} \textbf{15}, 654 (2024). DOI: 10.1038/s41467-024-44824-z.
\bibitem{ref40} Pfohl, S. R., Cole-Lewis, H., Sayres, R. \textit{et al}. A toolbox for surfacing health equity harms and biases in large language models (EquityMedQA). \textit{Nature Medicine} (2024). DOI: 10.1038/s41591-024-03258-2.
\bibitem{ref41} Kraljevic, Z., Bean, D., Shek, A. \textit{et al}. Large language models for medical forecasting---Foresight 2. arXiv:2412.10848 (2024).
\bibitem{ref42} Lekadir, K., Frangi, A. F., Porras, A. R. \textit{et al}. FUTURE-AI: international consensus guideline for trustworthy and deployable artificial intelligence in healthcare. \textit{The BMJ} (2025). DOI: 10.1136/bmj-2024-081554.
\bibitem{ref43} Shanghai AI Laboratory OpenMEDLab. PULSE: Pretrained Unified Large-scale Chinese Medical LM. GitHub release (2023). \url{https://github.com/openmedlab/PULSE}.
\bibitem{ref44} Wang, X., Chen, G. H., Song, D. \textit{et al}. CMB: A Comprehensive Medical Benchmark in Chinese. \textit{Proc. NAACL 2024 (Main)}. arXiv:2308.08833.
\bibitem{ref45} Jin, Q., Chen, F., Zhou, Y. \textit{et al}. Hidden flaws behind expert-level accuracy of multimodal GPT-4 vision in medicine. \textit{npj Digital Medicine} \textbf{7}, 190 (2024). DOI: 10.1038/s41746-024-01185-7.

\end{thebibliography}


\end{document}
